\documentclass[12pt, a4paper]{exam}
\usepackage{fontspec}
\usepackage{xunicode}
\usepackage{xltxtra}
\usepackage{polyglossia}
\usepackage{graphicx}
\usepackage{tikz}
\usetikzlibrary{positioning, calc}
\usepackage{geometry}
\usepackage{enumitem}
\usepackage{multicol}
\usepackage{array}
\usepackage{lastpage}
\usepackage{amssymb}
\usepackage{multirow}
\usepackage{amsmath}

\setmainlanguage{thai}
\XeTeXlinebreaklocale "th"
\XeTeXlinebreakskip = 0pt plus 1pt

% Page Margins
\geometry{top=0.8in, bottom=1in, left=0.8in, right=0.8in}

% Fonts
\newfontfamily\thaifont[
  Scale=1.33,
  Script=Thai,
  Path = ../fonts/,
  Extension = .ttf,
  BoldFont = *-Bold,
  ItalicFont = *-Italic,
  BoldItalicFont = *-BoldItalic
]{THSarabunNew}

\setmainfont[
  Scale=1.33,
  Path = ../fonts/,
  Extension = .ttf,
  BoldFont = *-Bold,
  ItalicFont = *-Italic,
  BoldItalicFont = *-BoldItalic
]{THSarabunNew}

\begin{document}

\begin{questions}

% Question 1
\question (25 คะแนน) จงตอบคำถามต่อไปนี้
\begin{enumerate}[label=\textbf{1.\arabic*)}]
    \item ให้ $A = \begin{bmatrix} 2 & -1 & 5 \\ 3 & 0 & -2 \\ 1 & 4 & -3 \end{bmatrix}$ เป็นเมทริกซ์มาตรฐาน (standard matrix) สำหรับการแปลงเชิงเส้น (linear transformation) $T : \mathbb{R}^3 \to \mathbb{R}^3$ จงหาค่าของ $T\left(\begin{bmatrix} 1 \\ -2 \\ 3 \end{bmatrix}\right)$
    \item กำหนดให้เมทริกซ์ $B = \begin{bmatrix} 4 & 1 & 0 & 0 \\ 0 & 4 & 0 & 0 \\ 0 & 0 & 2 & 1 \\ 0 & 0 & 0 & 2 \end{bmatrix}$ สามารถเขียนในรูปแบบของจอร์แดน (Jordan form) เป็น $\begin{bmatrix} M_1 & 0 \\ 0 & M_2 \end{bmatrix}$ จงระบุเมทริกซ์ $M_1$ และ $M_2$ สำหรับเมทริกซ์นี้
    \item จงหาเวกเตอร์เจาะจงทั่วไป ซึ่งมีลำดับขั้นเท่ากับ 2 (generalized eigenvector of rank 2) ที่สัมพันธ์กับค่าเจาะจง (eigenvalue) $\lambda = 3$ ของเมทริกซ์ $C = \begin{bmatrix} 4 & -1 \\ 1 & 2 \end{bmatrix}$
    \item กำหนดเมทริกซ์ $D = \begin{bmatrix} 7 & -4 \\ 5 & -2 \end{bmatrix}$ จงหาสมการลักษณะเฉพาะ (characteristic equation) และค่าเจาะจง (eigenvalues) ทั้งหมดของเมทริกซ์ $D$
\end{enumerate}
\vspace{\stretch{1}}
\newpage

% Question 2
\question (20 คะแนน) กําหนดให้ $T : \mathbb{R}^4 \to \mathbb{R}^3$ นิยามโดย
\[ T\left(\begin{bmatrix} x_1 \\ x_2 \\ x_3 \\ x_4 \end{bmatrix}\right) = \begin{bmatrix} x_1 - x_2 + x_3 + x_4 \\ x_1 + 2x_2 - x_3 + x_4 \\ 2x_1 + x_2 + 2x_4 \end{bmatrix} \]
\begin{enumerate}[label=\textbf{2.\arabic*)}]
    \item จงแสดงว่า $T$ เป็นการแปลงเชิงเส้น (Show that $T$ is a linear transformation.)
    \item จงหาฐานหลัก (a basis) ของ $Ker(T)$ และค่าของ $nullity(T)$
    \item จงหาฐานหลัก (a basis) ของ $Range(T)$ และค่าของ $rank(T)$
    \item จงตรวจสอบความถูกต้องโดยใช้ทฤษฎีบท Rank-Nullity และสรุปว่า $T$ เป็นฟังก์ชันหนึ่งต่อหนึ่ง (one-to-one) \mbox{หรือไม่} และเป็นฟังก์ชันทั่วถึง (onto) \mbox{หรือไม่}
\end{enumerate}
\vspace{\stretch{1}}
\newpage

% Question 3
\question (20 คะแนน) จงพิจารณาว่า เมทริกซ์ $A$ ที่กำหนดให้ต่อไปนี้เป็นเมทริกซ์ที่แปลงให้เป็นเมทริกซ์\mbox{ทแยงมุม}ได้\mbox{หรือไม่} (Determine if $A$ is diagonalizable.)
\[ A = \begin{bmatrix} 3 & -1 & 0 \\ -2 & 4 & 0 \\ 1 & 2 & 1 \end{bmatrix} \]
\begin{enumerate}[label=\textbf{3.\arabic*)}]
    \item จงหาค่าเจาะจง (eigenvalues) ทั้งหมดของ $A$ (คำใบ้: ค่าเจาะจงเป็นจำนวนเต็มบวกที่แตกต่างกันทั้งหมด)
    \item จงหาฐานหลักของปริภูมิเจาะจง (basis for eigenspace) สำหรับแต่ละค่าเจาะจง
    \item หาก $A$ เป็นเมทริกซ์ที่แปลงให้เป็นเมทริกซ์\mbox{ทแยงมุม}ได้ จงทำการหาเมทริกซ์ $P$ และเมทริกซ์\mbox{ทแยงมุม} (diagonal matrix) $D$ ซึ่งทำให้ $A = PDP^{-1}$ (ถ้าแปลงไม่ได้ให้ระบุเหตุผล)
    \item จงแสดงวิธีการใช้เมทริกซ์ $P$ และ $D$ นี้ในการคำนวณหาเมทริกซ์ $A^5$ (ตอบในรูปผลคูณเมทริกซ์ ไม่ต้องกระจายตัวเลข)
\end{enumerate}
\vspace{\stretch{1}}
\newpage

% Question 4
\question (11 คะแนน) กำหนดให้
\[ Q(x,y) = 7x^2 + 13y^2 + 8xy \]
อยู่ในรูปแบบกำลังสอง (in the quadratic form)
\begin{enumerate}[label=\textbf{4.\arabic*)}]
    \item (2 คะแนน) จงเขียน $Q(x,y)$ ในรูปสมการเมทริกซ์ $\mathbf{x}^T A \mathbf{x}$ โดยที่ $A$ เป็นเมทริกซ์สมมาตร (symmetric matrix)
    \item (6 คะแนน) จงหาเมทริกซ์เชิงตั้งฉาก $P$ (orthogonal matrix) ที่ทำให้ $P^T A P = D$ เป็นเมทริกซ์\mbox{ทแยงมุม}
    \item (3 คะแนน) จงทำการเปลี่ยนตัวแปร $\mathbf{x} = P \mathbf{x}'$ เพื่อให้รูปแบบกำลังสองของตัวแปรใหม่ ไม่มีเทอมผลคูณไขว้ปรากฏอยู่ (contains no cross-product terms) และระบุว่าเป็นสมการของกราฟรูปใด
\end{enumerate}
\vspace{\stretch{1}}
\newpage

% Question 5
\question (7 คะแนน) โจทย์ประยุกต์: การหมุนภาพในระนาบ (Rotation in $\mathbb{R}^2$)

บทนำ (Introduction): ในงานคอมพิวเตอร์กราฟิก การหมุนวัตถุรอบจุดกำเนิด (Rotation) เป็นการแปลงเชิงเส้นที่สำคัญ หากเราต้องการหมุนจุด $\mathbf{x} = \begin{bmatrix} x \\ y \end{bmatrix}$ ทวนเข็มนาฬิกาเป็นมุม $\theta$ ไปยังจุดใหม่ $\mathbf{x}' = \begin{bmatrix} x' \\ y' \end{bmatrix}$ ความสัมพันธ์นี้สามารถเขียนในรูปสมการเมทริกซ์ได้เป็น
\[ \mathbf{x}' = R_\theta \mathbf{x} \]
โดยที่ $R_\theta$ คือ Rotation Matrix:
\[ R_\theta = \begin{bmatrix} \cos\theta & -\sin\theta \\ \sin\theta & \cos\theta \end{bmatrix} \]
ซึ่งได้มาจากการพิจารณาภาพของเวกเตอร์ฐานหลัก $\mathbf{e}_1, \mathbf{e}_2$ ภายใต้การหมุน

จงตอบคำถามต่อไปนี้:
\begin{enumerate}[label=\textbf{5.\arabic*)}]
    \item \textbf{การคำนวณการหมุน (Forward Calculation) (2 คะแนน)} กำหนดให้มุมการหมุน $\theta = \frac{\pi}{6}$ ($30^\circ$) จงเขียนสมการเมทริกซ์แทนการแปลงนี้ และจงหาพิกัดใหม่ $\mathbf{x}'$ ของจุด $\mathbf{x} = \begin{bmatrix} 2 \\ 4 \end{bmatrix}$ หลังจากถูกหมุนทวนเข็มนาฬิกาด้วยมุมนี้ (กำหนดให้ $\cos 30^\circ = \frac{\sqrt{3}}{2}, \sin 30^\circ = \frac{1}{2}$)
    \item \textbf{การหาค่ามุม (Inverse Problem) (2 คะแนน)} สมมติว่าโปรแกรมกราฟิกทำการหมุนจุด $\mathbf{x} = \begin{bmatrix} 1 \\ 0 \end{bmatrix}$ ไปยังจุดใหม่ $\mathbf{x}' = \begin{bmatrix} 0 \\ 1 \end{bmatrix}$ โดยการหมุนรอบจุดกำเนิดในทิศทางทวนเข็มนาฬิกา จงเขียนสมการเมทริกซ์
    \[ \begin{bmatrix} 0 \\ 1 \end{bmatrix} = \begin{bmatrix} \cos\theta & -\sin\theta \\ \sin\theta & \cos\theta \end{bmatrix} \begin{bmatrix} 1 \\ 0 \end{bmatrix} \]
    เพื่อหาค่า มุมของการหมุน $\theta$ (ในหน่วยเรเดียน)
    \item \textbf{การพิสูจน์ (Theoretical Proof) (3 คะแนน)} จงพิสูจน์สมบัติการหมุนซ้ำ $n$ ครั้ง โดยเขียนในรูปการคูณเมทริกซ์ว่า
    \[ (R_\theta)^n = \begin{bmatrix} \cos(n\theta) & -\sin(n\theta) \\ \sin(n\theta) & \cos(n\theta) \end{bmatrix} \]
    สำหรับจำนวนเต็มบวก $n$ ใด ๆ
\end{enumerate}

\end{questions}
\end{document}
